\section{AMC 12 12B 2010}
        \begin{enumerate}
        	\item (\textit{AMC 12 12B 2010}) Makarla attended two meetings during her $9$-hour work day. The first meeting took $45$ minutes and the second meeting took twice as long. What percent of her work day was spent attending meetings?


 $\textbf{(A)}\ 15 \qquad \textbf{(B)}\ 20 \qquad \textbf{(C)}\ 25 \qquad \textbf{(D)}\ 30 \qquad \textbf{(E)}\ 35$

 \textbf{Answer: }C



\textbf{Solution}

The total number of minutes in her $9$-hour work day is $9 \times 60 = 540.$The total amount of time spend in meetings in minutes is $45 + 45 \times 2 = 135.$The answer is then $\frac{135}{540}$ $= \boxed{25\%}$ or  $\boxed{(C)}$



	\item (\textit{AMC 12 12B 2010}) A big $L$ is formed as shown. What is its area?


 
                    \begin{figure}[H]
                        \centering
                        \includegraphics[width=0.5\textwidth]{Problem/AMC_12_12B_2010_2_Problem_Figure_1.png}
                    \end{figure}
                    
$\textbf{(A)}\ 22 \qquad \textbf{(B)}\ 24 \qquad \textbf{(C)}\ 26 \qquad \textbf{(D)}\ 28 \qquad \textbf{(E)}\ 30$

 \textbf{Answer: }A



\textbf{Solution}

We find the area of the big rectangle to be $8 \times 5 = 40$, and the area of the smaller rectangle to be $(8 - 2) \times (5 - 2) = 18$. The answer is then $40 - 18 = 22$ $(A)$.

 


 


	\item (\textit{AMC 12 12B 2010}) A ticket to a school play cost $x$ dollars, where $x$ is a whole number. A group of 9<sup>th</sup> graders buys tickets costing a total of $$48$, and a group of 10<sup>th</sup> graders buys tickets costing a total of $$64$. How many values for $x$ are possible?


 $\textbf{(A)}\ 1 \qquad \textbf{(B)}\ 2 \qquad \textbf{(C)}\ 3 \qquad \textbf{(D)}\ 4 \qquad \textbf{(E)}\ 5$

 \textbf{Answer: }E



\textbf{Solution}

We find the greatest common factor of $48$ and $64$ to be $16$. The number of factors of $16$ is $5$ which is the answer $(E)$.

 



\textbf{Solution 2}

Since the price of one ticket is a whole number, it must be a divisor of both $48$ and $64$. An integer is a common divisor of two numbers if and only if it is divisible by their difference. Hence, all the possible prices are positive divisors of $16$, or $1$, $2$, $4$, $8$, and $16$, making the answer $5$, or E. - By Monkey\_King

 


	\item (\textit{AMC 12 12B 2010}) A month with $31$ days has the same number of Mondays and Wednesdays. How many of the seven days of the week could be the first day of this month?


 $\textbf{(A)}\ 2 \qquad \textbf{(B)}\ 3 \qquad \textbf{(C)}\ 4 \qquad \textbf{(D)}\ 5 \qquad \textbf{(E)}\ 6$

 \textbf{Answer: }B



\textbf{Solution}

$31 \equiv 3 \pmod {7}$ so the week cannot start with Saturday, Sunday, Tuesday or Wednesday as that would result in an unequal number of Mondays and Wednesdays. Therefore, Monday, Thursday, and Friday are valid so the answer is $\boxed{B}$.

 


	\item (\textit{AMC 12 12B 2010}) Lucky Larry's teacher asked him to substitute numbers for $a$, $b$, $c$, $d$, and $e$ in the expression $a-(b-(c-(d+e)))$ and evaluate the result. Larry ignored the parentheses but added and subtracted correctly and obtained the correct result by coincidence. The numbers Larry substituted for $a$, $b$, $c$, and $d$ were $1$, $2$, $3$, and $4$, respectively. What number did Larry substitute for $e$?


 $\textbf{(A)}\ -5 \qquad \textbf{(B)}\ -3 \qquad \textbf{(C)}\ 0 \qquad \textbf{(D)}\ 3 \qquad \textbf{(E)}\ 5$

 \textbf{Answer: }D



\textbf{Solution}

We simply plug in the numbers
 \[1 - 2 - 3 - 4 + e = 1 - (2 - (3 - (4 + e)))\]
 \[-8 + e = -2 - e\]
 \[2e = 6\]
 \[e = 3 \;\;(D)\]



	\item (\textit{AMC 12 12B 2010}) At the beginning of the school year, $50\%$ of all students in Mr. Wells' math class answered "Yes" to the question "Do you love math", and $50\%$ answered "No." At the end of the school year, $70\%$ answered "Yes" and $30\%$ answered "No." Altogether, $x\%$ of the students gave a different answer at the beginning and end of the school year. What is the difference between the maximum and the minimum possible values of $x$?


 $\textbf{(A)}\ 0 \qquad \textbf{(B)}\ 20 \qquad \textbf{(C)}\ 40 \qquad \textbf{(D)}\ 60 \qquad \textbf{(E)}\ 80$

 \textbf{Answer: }D



\textbf{Solution}

The minimum possible value would be $70 - 50 = 20\%$. The maximum possible value would be $30 + 50 = 80\%$. The difference is $80 - 20 = \boxed{\textbf{(D) }60}$.

 


	\item (\textit{AMC 12 12B 2010}) Shelby drives her scooter at a speed of $30$ miles per hour if it is not raining, and $20$ miles per hour if it is raining. Today she drove in the sun in the morning and in the rain in the evening, for a total of $16$ miles in $40$ minutes. How many minutes did she drive in the rain?


 $\textbf{(A)}\ 18 \qquad \textbf{(B)}\ 21 \qquad \textbf{(C)}\ 24 \qquad \textbf{(D)}\ 27 \qquad \textbf{(E)}\ 30$

 \textbf{Answer: }C



\textbf{Solution 1}

Let $x$ be the time it is not raining, and $y$ be the time it is raining, in hours.

 We have the system: $30x+20y=16$ and $x+y=\frac{2}{3}$

 Solving gives $x=\frac{4}{15}$ and $y=\frac{2}{5}$

 We want $y$ in minutes, $\frac{2}{5} \cdot 60=24 \Rightarrow C$




\textbf{Solution 2}

Let $x$ be the time she drove in the rain. Thus, the number of minutes she did not drive in the rain is $40-x$ . 

 Since we are calculating the time in minutes, it is best to convert the speeds in minutes. Thus, the speed per minute when it is not raining is $\frac{1}{2}$ miles per minute, and $\frac{1}{3}$ miles per minute when it is raining. Thus, we have the equation, $\frac{1}{2} \cdot (40-x) + \frac{1}{3} \cdot x = 16$ 

 Solving, gives $x$ = $24$ , so the amount of time she drove in the rain is $24$ minutes.

 


	\item (\textit{AMC 12 12B 2010}) Every high school in the city of Euclid sent a team of $3$ students to a math contest. Each participant in the contest received a different score. Andrea's score was the median among all students, and hers was the highest score on her team. Andrea's teammates Beth and Carla placed $37$<sup>th</sup> and $64$<sup>th</sup>, respectively. How many schools are in the city?


 $\textbf{(A)}\ 22 \qquad \textbf{(B)}\ 23 \qquad \textbf{(C)}\ 24 \qquad \textbf{(D)}\ 25 \qquad \textbf{(E)}\ 26$

 \textbf{Answer: }B



\textbf{Solution 1}

There are $x$ schools. This means that there are $3x$ people. Because no one's score was the same as another person's score, that means that there could only have been $1$ median score. This implies that $x$ is an odd number. $x$ cannot be less than $23$, because there wouldn't be a $64$th place if $x$ was. $x$ cannot be greater than $23$ either, because that would tie Andrea and Beth or Andrea's place would be worse than Beth's. Thus, the only possible answer is $\boxed{\mathbf{(B)}\ 23}$.

 



\textbf{Solution 2}

Let $a$ be Andrea's place. We know that she was the highest on her team, so $a < 37$. 

 Since $a$ is the median, there are $a-1$ to the left and right of the median, so the total number of people is $2a-1$ and the number of schools is $(2a-1)/3$. This implies that $2a-1 \equiv 0 \pmod{3} \implies a \equiv 2 \pmod{3}$. 

 Also, since $2a-1$ is the rank of the last-place person, and one of Andrea's teammates already got 64th place, $2a-1 > 64 \implies a \ge 33$.

 Putting it all together: $33 \le a < 37$ and $a \equiv 2 \pmod{3}$, so clearly $a = 35$, and the number of schools as we got before is $(2a-1)/3 = 69/3 = \boxed{\mathbf{(B)}\ 23}$.

 



\textbf{Solution 3}

Since all students have distinct scores, we can assume the scores are $1$ through $S$, where $S$ is the total number of students. From Solution 1, we know that $S$ is odd.

 Andrea’s score is the median, so$A = \frac{S+1}{2}$.Beth and Carla have scores$B = S - 37$, $\quad C = S - 64$.Since Andrea has the highest score among her teammates, we have$A > B \quad \text{and} \quad A > C$.Thus, $\frac{S+1}{2} > S - 37 \quad \text{and} \quad \frac{S+1}{2} > S - 64$.

 Solving these inequalities:$\frac{S+1}{2} > S - 37 \Rightarrow S < 75$,$\frac{S+1}{2} > S - 64 \Rightarrow S < 129$.Hence, $S < 75$.We also have : $B + C > A$.Substituting:$(S - 37) + (S - 64) > \frac{S+1}{2}$.This simplifies to$2S - 101 > \frac{S+1}{2}$$\Rightarrow 4S - 202 > S + 1$$\Rightarrow 3S > 203$$\Rightarrow S > \frac{203}{3}$.

 Since students are grouped into teams of 3, let$S = 3t$.Then$\frac{203}{3} < S < 75$.$\Rightarrow 22 < t < 25$Because t has to be an odd number in this range we get :$t = 23$.

 $\boxed{\mathbf{(B)\ 23}}$



	\item (\textit{AMC 12 12B 2010}) Let $n$ be the smallest positive integer such that $n$ is divisible by $20$, $n^2$ is a perfect cube, and $n^3$ is a perfect square. What is the number of digits of $n$?


 $\textbf{(A)}\ 3 \qquad \textbf{(B)}\ 4 \qquad \textbf{(C)}\ 5 \qquad \textbf{(D)}\ 6 \qquad \textbf{(E)}\ 7$

 \textbf{Answer: }E



\textbf{Solution}

We know that $n^2 = k^3$ and $n^3 = m^2$. Cubing and squaring the equalities respectively gives $n^6 = k^9 = m^4$. Let $a = n^6$. Now we know $a$ must be a perfect $36$-th power because $lcm(9,4) = 36$, which means that $n$ must be a perfect $6$-th power. The smallest number whose sixth power is a multiple of $20$ is $10$, because the only prime factors of $20$ are $2$ and $5$, and $10 = 2 \times 5$. Therefore our is equal to number $10^6 = 1000000$, with $7$ digits $\Rightarrow \boxed {E}$.

 



\textbf{Solution 2 (Chinese Remainder Theorem)}

Let $n = 2^2 \cdot 5 \cdot x$ for some integer $x$, then we know that \begin{align*}n^2 = 2^4 \cdot 5^2 \cdot x^2 = m_0^3\\n^3 = 2^6 \cdot 5^3 \cdot x^3 = m_1^2\end{align*}
Since we have $2^4 \cdot 5^2 \cdot x^2$ a power of 3, that means, by the Fundamental Theorem of Arithmetic, the exponents of the primes must each be divisible by 3. This also means that there's no point in $p \mid x$ if $p \not \in \{2,3\}$, since that would just serve to increase $n$ and we want $n$ to be minimal.

 This means we can write $x = 2^{\alpha} \cdot 5^{\beta}$ and thus we have $2^4 \cdot 5^2 \cdot 2^{2\alpha} \cdot 5^{2\beta} = 2^{4 + 2\alpha} \cdot 5^{2 + 2\beta} = m_0^3$, therefore we know \begin{align*}4 + 2\alpha \equiv 0 \pmod{3}\\2 + 2\beta \equiv 0 \pmod{3}\end{align*}
and doing the same with the third equation being a square, we have\begin{align*}6 + 3\alpha \equiv 0 \pmod{2}\\3 + 3\beta \equiv 0 \pmod{2}\end{align*}
We can shift this around into two systems involving the same variables, as follows:\begin{align*}\alpha \equiv 0 \pmod{2}\\1 + 2\alpha \equiv 0 \pmod{3}\end{align*}
and\begin{align*}1 + \beta \equiv 0 \pmod{2} \\2 + 2\beta \equiv 0 \pmod{3}\\\end{align*}
And using \href{https://artofproblemsolving.com/wiki/index.php?title=Chinese\_Remainder\_Theorem}{Chinese Remainder Theorem} to solve this, we get $\alpha = 4$ and $\beta = 5$, and plugging that back in yields $n = 2^6 \cdot 5^6$.

 Thus, our answer is $\lfloor{ \log_{10}(2^6 \cdot 5^6) + 1 \rfloor} \Rightarrow \boxed {E}$. 

 


	\item (\textit{AMC 12 12B 2010}) The average of the numbers $1, 2, 3,\cdots, 98, 99,$ and $x$ is $100x$. What is $x$?


 $\textbf{(A)}\ \dfrac{49}{101} \qquad \textbf{(B)}\ \dfrac{50}{101} \qquad \textbf{(C)}\ \dfrac{1}{2} \qquad \textbf{(D)}\ \dfrac{51}{101} \qquad \textbf{(E)}\ \dfrac{50}{99}$

 \textbf{Answer: }B



\textbf{Solution}

We first sum the first $99$ numbers: $\frac{99(100)}{2}=99\cdot50$. Then, we know that the sum of the series is $99\cdot50+x$. There are $100$ terms, so we can divide this sum by $100$ and set it equal to $100x$:
 \[\frac{99\cdot50+x}{100}=100x \Rightarrow 99\cdot50=100^2x-x\Rightarrow \frac{99\cdot50}{100^2-1}=x\]Using difference of squares:
 \[x=\frac{99\cdot50}{101\cdot99}=\frac{50}{101}\]Thus, the answer is $\boxed{\textbf{(B) }\frac{50}{101}}$.

 


	\item (\textit{AMC 12 12B 2010}) A palindrome between $1000$ and $10,000$ is chosen at random. What is the probability that it is divisible by $7$?


 $\textbf{(A)}\ \dfrac{1}{10} \qquad \textbf{(B)}\ \dfrac{1}{9} \qquad \textbf{(C)}\ \dfrac{1}{7} \qquad \textbf{(D)}\ \dfrac{1}{6} \qquad \textbf{(E)}\ \dfrac{1}{5}$

 \textbf{Answer: }E



\textbf{Solution 1}

View the palindrome as some number with form (decimal representation):$a_3 \cdot 10^3 + a_2 \cdot 10^2 + a_1 \cdot 10 + a_0$. But because the number is a palindrome, $a_3 = a_0, a_2 = a_1$. Recombining this yields $1001a_3 + 110a_2$. 1001 is divisible by 7, which means that as long as $a_2 = 0$, the palindrome will be divisible by 7. This yields 9 palindromes out of 90 ($9 \cdot 10$) possibilities for palindromes. However, if $a_2 = 7$, then this gives another case in which the palindrome is divisible by 7. This adds another 9 palindromes to the list, bringing our total to $18/90 = \boxed {\frac{1}{5} } = \boxed {E}$




\textbf{Solution (Divisibility Rules)}

We can notice the palindrome is of the form $\overline{abba}$. Then, by the divisibility rule of $7$, $7$ must divide 
 \[100a+11b-2a = 98a+11b.\] This nicely simplifies to the fact that $7 \mid 4b,$ so $b$ is clearly $0$ or $7$. This gives us $9 \cdot 2$ total choices for the palindrome divisible by $7$, divided by $9 \cdot 10$ total choices for $\overline{abba}$, giving us an answer of $\boxed{\text{(E)}} \ \dfrac{1}{5}$.

 \textbf{Addendum (Alternate)} $7\mid 1001a^3+110b^2$ and $1001 \equiv 0 \pmod 7$. Knowing that $a$ does not factor (pun intended) into the problem, note 110's prime factorization and $7\mid b$. There are only 10 possible digits for $b$, 0 through 9, but $7\mid b$ only holds if $b=0, 7$. This is 2 of the 10 digits, so $\frac{2}{10}=\boxed{\textbf{(E) }\frac{1}{5}}$




\textbf{Solution 4}

The palindrome is in the form of $\overline{abba}$. There are $9$ possible values of $a$ and 10 possible values for $b$. Thus, there are $90$ four-digit palindromes. The smallest palindrome is $1001$, which is a multiple of $7$. Note how each palindrome with the same thousands and ones digit increase by $110$ (i.e., $1111-1001=110$, $1221-1111=110$, etc). $110 \equiv 5$ (mod $7$). Since $5$ and $7$ is coprime, to get another palindrome that is a multiple of $7$, we need to increase by $770$. $1001+770=1771$. Next, when the thousands place is $2$, we start with $2002$, which is a multiple of $1001$, so $2002 \equiv 0$ (mod $7$). Therefore, $2772$ is also a satisfactory palindrome. Similarly, for every thousands place, there are $2$ satisfactory palindromes, which means there are a total of $2 \cdot 9 = 18$ satisfactory palindromes. Thus, the answer is $\frac{18}{90} = \frac{1}{5}$ or $\boxed{\text{(E)}\ \frac{1}{5}}$.

 


	\item (\textit{AMC 12 12B 2010}) For what value of $x$ does


 
 \[\log_{\sqrt{2}}\sqrt{x}+\log_{2}{x}+\log_{4}{x^2}+\log_{8}{x^3}+\log_{16}{x^4}=40?\]
$\textbf{(A)}\ 8 \qquad \textbf{(B)}\ 16 \qquad \textbf{(C)}\ 32 \qquad \textbf{(D)}\ 256 \qquad \textbf{(E)}\ 1024$

 \textbf{Answer: }D



\textbf{Solution 1}


 \[\log_{\sqrt{2}}\sqrt{x} + \log_2x + \log_4(x^2) + \log_8(x^3) + \log_{16}(x^4)  = 40\]

 \[\frac{1}{2} \frac{\log_2x}{\log_2\sqrt{2}} + \log_2x + \frac{2\log_2x}{\log_24} + \frac{3\log_2x}{\log_28} + \frac{4\log_2x}{\log_216}  = 40\]

 \[\log_2x + \log_2x + \log_2x + \log_2x + \log_2x = 40\]

 \[5\log_2x = 40\]

 \[\log_2x = 8\]

 \[x = 256 \;\; (D)\]




\textbf{Solution 2}

Using the fact that $\log_{x^n}{y^n} = \log_{x}{y}$, we see that the equation becomes $\log_{2}{x} + \log_{2}{x} + \log_{2}{x} + \log_{2}{x} + \log_{2}{x} = 40$. Thus, $5\log_{2}{x} = 40$ and $\log_{2}{x} = 8$, so $x = 2^8 = 256$, or $\boxed{(D)}$.

 


 


	\item (\textit{AMC 12 12B 2010}) In $\triangle ABC$, $\cos(2A-B)+\sin(A+B)=2$ and $AB=4$. What is $BC$?


 $\textbf{(A)}\ \sqrt{2} \qquad \textbf{(B)}\ \sqrt{3} \qquad \textbf{(C)}\ 2 \qquad \textbf{(D)}\ 2\sqrt{2} \qquad \textbf{(E)}\ 2\sqrt{3}$

 \textbf{Answer: }C



\textbf{Solution}

We note that $-1$ $\le$ $\sin x$ $\le$ $1$ and $-1$ $\le$ $\cos x$ $\le$ $1$. Therefore, there is no other way to satisfy this equation other than making both $\cos(2A-B)=1$ and $\sin(A+B)=1$, since any other way would cause one of these values to become greater than 1, which contradicts our previous statement.From this we can easily conclude that $2A-B=0^{\circ}$ and $A+B=90^{\circ}$ and solving this system gives us $A=30^{\circ}$ and $B=60^{\circ}$. It is clear that $\triangle ABC$ is a $30^{\circ},60^{\circ},90^{\circ}$ triangle with $BC=2$　$\Longrightarrow$　$(C)$.

 


	\item (\textit{AMC 12 12B 2010}) Let $a$, $b$, $c$, $d$, and $e$ be positive integers with $a+b+c+d+e=2010$ and let $M$ be the largest of the sums $a+b$, $b+c$, $c+d$ and $d+e$. What is the smallest possible value of $M$?


 $\textbf{(A)}\ 670 \qquad \textbf{(B)}\ 671 \qquad \textbf{(C)}\ 802 \qquad \textbf{(D)}\ 803 \qquad \textbf{(E)}\ 804$

 \textbf{Answer: }B



\textbf{Solution 1}

We want to try make $a+b$, $b+c$, $c+d$, and $d+e$ as close as possible so that $M$, the maximum of these, is smallest.

 Notice that $2010=670+670+670$. In order to express $2010$ as a sum of $5$ numbers, we must split up some of these numbers. There are two ways to do this (while keeping the sum of two numbers as close as possible): $2010=670+1+670+1+668$ or $2010=670+1+669+1+669$. We see that in both cases, the value of $M$ is $671$, so the answer is $671 \Rightarrow \boxed{B}$.

 



\textbf{Solution 2}

Since $a + b \le M$, $d + e \le M$, and $c < b + c \le M$, we have that $2010 = a + b + c + d + e < 3M$. Hence, $M > 670$, or $M \ge 671$.

 For the values $(a,b,c,d,e) = (669,1,670,1,669)$, $M = 671$, so the smallest possible value of $M$ is $\boxed{671}$. The answer is (B).

 



\textbf{Solution 3}

Notice that only the sums of adjacent numbers matter. (For example, a & c could be extremely high, as long as b is relatively low.)Therefore creating "mountains" and "valleys" is the best way to lower the sum of adjacent numbers. We can do

 
1. (high, low, high, low, high)

 
or

 
2. (low, high, low, high, low)

 
In the extreme case that each "low" = $0$, $2010$ will be divided into either $3$ or $2$ numbers for cases 1 and 2, respectively. Obviously dividing by $3$ will yield a lower number, so we consider case 1.

 Dividing $2010$ by $3$ yields $670$, or

 
($670, 0, 670, 0, 670$)

 
However, all five numbers must be positive, and the closest we can get to this is

 
($668, 3, 668, 3, 668$)

 
The lowest possible sum of two adjacent numbers then becomes $671$, or $\boxed{B}$.

 


	\item (\textit{AMC 12 12B 2010}) For how many ordered triples $(x,y,z)$ of nonnegative integers less than $20$ are there exactly two distinct elements in the set $\{i^x, (1+i)^y, z\}$, where $i=\sqrt{-1}$?


 $\textbf{(A)}\ 149 \qquad \textbf{(B)}\ 205 \qquad \textbf{(C)}\ 215 \qquad \textbf{(D)}\ 225 \qquad \textbf{(E)}\ 235$

 \textbf{Answer: }D



\textbf{Solution}

We have either $i^{x}=(1+i)^{y}\neq z$, $i^{x}=z\neq(1+i)^{y}$, or $(1+i)^{y}=z\neq i^x$.

 $i^{x}=(1+i)^{y}$ only occurs when it is $1$. $(1+i)^{y}=1$ has only one solution, namely, $y=0$. $i^{x}=1$ has five solutions between 0 and 19, $x=0, x=4, x=8, x=12$, and $x=16$. $z\neq 1$ has nineteen integer solutions between zero and nineteen. So for $i^{x}=(1+i)^{y}\neq z$, we have $5\cdot 1\cdot 19=95$ ordered triples.

 For $i^{x}=z\neq(1+i)^{y}$, again this only occurs at $1$. $(1+i)^{y}\neq 1$ has nineteen solutions, $i^{x}=1$ has five solutions, and $z=1$ has one solution, so again we have $5\cdot 1\cdot 19=95$ ordered triples.

 For $(1+i)^{y}=z\neq i^x$, this occurs at $1$ and $16$. $(1+i)^{y}=1$ and $z=1$ both have one solution while $i^{x}\neq 1$ has fifteen solutions. $(1+i)^{y}=16$ and $z=16$ both have one solution, namely, $y=8$ and $z=16$, while $i^{x}\neq 16$ has twenty solutions ($i^x$ only cycles as $1, i, -1, -i$). So we have $15\cdot 1\cdot 1+20\cdot 1\cdot 1=35$ ordered triples.

 In total we have ${95+95+35=\boxed{\text{(D) }225}}$ ordered triples

 \textbf{Small Clarification} To more clearly see why the reasoning above is true, try converting the complex numbers into exponential form. That way, we can more easily raise the numbers to $x$, $y$ and $z$ respectively.

 


	\item (\textit{AMC 12 12B 2010}) Positive integers $a$, $b$, and $c$ are randomly and independently selected with replacement from the set $\{1, 2, 3,\dots, 2010\}$. What is the probability that $abc + ab + a$ is divisible by $3$?


 $\textbf{(A)}\ \dfrac{1}{3} \qquad \textbf{(B)}\ \dfrac{29}{81} \qquad \textbf{(C)}\ \dfrac{31}{81} \qquad \textbf{(D)}\ \dfrac{11}{27} \qquad \textbf{(E)}\ \dfrac{13}{27}$

 \textbf{Answer: }E



\textbf{Solution 1}

We group this into groups of $3$, because $3|2010$. This means that every residue class mod 3 has an equal probability.

 If $3|a$, we are done. There is a probability of $\frac{1}{3}$ that that happens.

 Otherwise, we have $3|bc+b+1$, which means that $b(c+1) \equiv 2\pmod{3}$. So either 
 \[b \equiv 1 \pmod{3}, c \equiv 1 \pmod{3}\] or 
 \[b \equiv 2 \pmod {3}, c \equiv 0 \pmod 3\] which will lead to the property being true. There is a $\frac{1}{3}\cdot\frac{1}{3}=\frac{1}{9}$ chance for each bundle of cases to be true. Thus, the total for the cases is $\frac{2}{9}$. But we have to multiply by $\frac{2}{3}$ because this case ($3|bc+b+1$ ) only needs to happen when $a$ isn’t divisible by 3,  which has a $\frac{2}{3}$ chance. So the total is actually $\frac{4}{27}$.

 The grand total is 
 \[\frac{1}{3} + \frac{4}{27} = \boxed{\text{(E) }\frac{13}{27}.}\]




\textbf{Solution 2 (Minor change from Solution 1)}

Just like solution 1, we see that there is a $\frac{1}{3}$ chance of $3|a$ and $\frac{2}{9}$ chance of $3|1+b+bc$ 

 Now, we can just use PIE (Principles of Inclusion and Exclusion) to get our answer to be $\frac{1}{3}+\frac{2}{9}-\frac{1}{3}\cdot\frac{2}{9} = \boxed{\text{(E) } \frac{13}{27}}$

 -Conantwiz2023

 



\textbf{Solution 3 (Fancier version of Solution 1)}

As with solution one, we conclude that if $a\equiv0\mod 3$ then the requirements are satisfied. We then have:
 \[a(bc+c)+a\equiv0 \mod 3\]
 \[a(bc+c)\equiv-a \mod 3\]
 \[c(b+1)\equiv-1 \mod 3\]
 \[b+1\equiv \frac{-1}{c} \mod 3\]Which is true for one $b$ when $c\not\equiv 0 \mod 3$ because the integers$\mod 3$ form a field under multiplication and addition with absorbing element $0$.

 This gives us $P=\frac{1}{3}+\frac{4}{27}=\boxed{\text{(E) }\frac{13}{27}}$.

 


	\item (\textit{AMC 12 12B 2010}) The entries in a $3 \times 3$ array include all the digits from $1$ through $9$, arranged so that the entries in every row and column are in increasing order. How many such arrays are there?


 $\textbf{(A)}\ 18 \qquad \textbf{(B)}\ 24 \qquad \textbf{(C)}\ 36 \qquad \textbf{(D)}\ 42 \qquad \textbf{(E)}\ 60$

 \textbf{Answer: }D



\textbf{Solution 1}

Observe that all tables must have 1s and 9s in the corners, 8s and 2s next to those corner squares, and 4-6 in the middle square. Also note that for each table, there exists a valid table diagonally symmetrical across the diagonal extending from the top left to the bottom right. 

 


 
            \begin{itemize}
        \item \textbf{Case 1: Center 4}
            \end{itemize}
            
 \[\begin{tabular}{|c|c|c|} \hline 1&2&\\ \hline 3&4&8\\ \hline &&9\\ \hline \end{tabular} \;\;\; \begin{tabular}{|c|c|c|} \hline 1&2&\\ \hline 3&4&\\ \hline &8&9\\ \hline \end{tabular}\]

 3 necessarily must be placed as above. Any number could fill the isolated square, but the other 2 are then invariant. So, there are 3 cases each and 6 overall cases. Given diagonal symmetry, alternate 2 and 8 placements yield symmetrical cases. $2*6=12$

            \begin{itemize}
        \item \textbf{Case 2: Center 5}
            \end{itemize}
            
 \[\begin{tabular}{|c|c|c|} \hline 1&2&3\\ \hline 4&5&\\ \hline &8&9\\ \hline \end{tabular} \;\;\; \begin{tabular}{|c|c|c|} \hline 1&2&\\ \hline 3&5&\\ \hline &8&9\\ \hline \end{tabular} \;\;\; \begin{tabular}{|c|c|c|} \hline 1&2&\\ \hline 3&5&8\\ \hline &&9\\ \hline \end{tabular} \;\;\; \begin{tabular}{|c|c|c|} \hline 1&2&3\\ \hline 4&5&8\\ \hline &&9\\ \hline \end{tabular}\]

 Here, no 3s or 7s are assured, but this is only a teensy bit trickier and messier. WLOG, casework with 3 instead of 7 as above. Remembering that $4<5$, logically see that the numbers of cases are then 2,3,3,1 res

 


	\item (\textit{AMC 12 12B 2010}) A frog makes $3$ jumps, each exactly $1$ meter long. The directions of the jumps are chosen independently at random. What is the probability that the frog's final position is no more than $1$ meter from its starting position?


 $\textbf{(A)}\ \dfrac{1}{6} \qquad \textbf{(B)}\ \dfrac{1}{5} \qquad \textbf{(C)}\ \dfrac{1}{4} \qquad \textbf{(D)}\ \dfrac{1}{3} \qquad \textbf{(E)}\ \dfrac{1}{2}$

 \textbf{Answer: }C



\textbf{Solution 1 (Complex Numbers)}

We will let the moves be complex numbers $a$, $b$, and $c$, each of magnitude one. The  starts on the origin. It is relatively easy to show that exactly one element in the set
 \[\{|a + b + c|, |a + b - c|, |a - b + c|, |a - b - c|\}\]has magnitude less than or equal to $1$. (Can you show how?) Hence, the probability is $\boxed{\text{(C)} \frac {1}{4}}$.

 



\textbf{Solution 2 (Simple Calculus)}

Yes, we are pulling out calculus...

 
Represent every jump as a circle of radius 1. The first circle is a circle of radius 1 centered on the origin. WLOG, assume the first jump lands on $(1,0)$. Then, the third circle could be centered anywhere on the second circle, which is itself centered on $(1,0)$. Let us define $f(\theta)$ as the value of the length of the first circle that lies within the area of the third circle in terms of the angle formed by the two points of intersection and either circle's center (symmetry, you chose!). The intersection of the two circles should form a geometrical lens shape. By sectors,
 \[f(\theta)=\frac{\theta}{2\pi}\cdot2\pi=\theta\] As the angle or angle of intersection continuously decreases from $\pi$ (when the third circle is on top of the second circle) to $0$ (when the third circle is only touching the first circle at one spot $(1,0)$, I just need to find the average value of this function to find the average arc length where the third jump could land to satisfy the problem. To do this, I can apply average function value with our old buddy calculus,
 \[\frac{1}{\pi}\int_0^{\pi}\theta\;d\theta=\frac{\pi}{2}\]The probability that the third jump will land on this arc length is just the arc length divided by the circumference, or $\frac{\frac{\pi}{2}}{2\pi}=\frac{1}{4} \implies \boxed{C}$




\textbf{Solution 3 (Geometric)}

The first  hop doesn't matter because no matter where the  hops, it lands on the border of the circle you want it to end in.  The remaining places that the frog can jump to form a disk of radius 2 centered at the spot on which the  first landed, and every point in the disk of radius 2 is equally likely to be reached in two jumps.

 
No matter where we start, we will have the small circle tangent to a point on the big circle. This is just like how $\frac{1}{2}x^2$ and $x^2$ are tangent. The area ratio of the two circles is 
 \[\frac{\pi}{4\pi} = \boxed{\frac{1}{4} \text{(C)}}\]




\textbf{Solution 4 (Easier Geometric Casework)}

Firstly, we know that the first jump will land on a circle of radius one around the frog's starting position. Then, the frog has two possibilities for where he could jump next.

 $\textbf{Case 1: }$ The frog's next jump lands outside of the circle with radius 1. The frog's possible area he could land on is the area formed by a circle of radius 2, and the probability of him landing outside of the radius 1 circle into the radius 2 circle is 
 \[\frac{2^2 \cdot \pi - 1^2 \cdot \pi}{2^2 \cdot \pi} = \frac{3}{4}\].

 To find the probability of the frog landing back in the circle of radius 1 on the third jump, we need to find the range that the circle could have traveled. After the second jump in case 1, the circle's position range from still on the circumference of the radius 1 circle, to being on the circumference of the radius 2 circle.

 When on the radius 1 circle, the frog has, similarily, a $\frac{1}{4}$ chance of jumping back in, while on the radius 2 circle, the frog has 0 chance of jumping back in the radius 1 circle, so the range is $[0, \frac{1}{4}]$. Taking the average, we get $\frac{1}{8}$, and that the frog have a total of 
 \[\frac{3}{4} \cdot \frac{1}{8} = \frac{3}{32}\] percent chance to jump back in in this case.

 $\textbf{Case 2: }$ The frog's next jump lands outside of the circle with radius 1. The frog's possible area he could land on is the area formed by a circle of radius 2, and the probability of him landing inside of the radius 1 circle is simply 
 \[\frac{1^2 \cdot \pi}{2^2 \cdot \pi} = \frac{1}{4}\].

 Similarily, to find the probability of the frog landing back in the circle of radius 1 on the third jump, we need to find the range that the circle could have traveled. After the second jump in case 2, the circle's position range from still on the circumference of the radius 1 circle, to being back to its starting point.

 When on the radius 1 circle, the frog has, a $\frac{1}{4}$ chance of jumping back in, while when the frog is on its starting point, the frog is certain to still be within the circle after one jump. Therefore, our range is going to be $[\frac{1}{4}, 1]$. Taking the average, we get $\frac{5}{8}$, and that the frog have a total of 
 \[\frac{1}{4} \cdot \frac{5}{8} = \frac{5}{32}\] percent chance to jump back in in this case. 

 Adding the probability of the two cases, we get that the frog has a 
 \[\frac{5}{32} + \frac{3}{32} = \frac{8}{32} = \frac{1}{4}\] chance of jumping back into no more than $1$ meter from its starting position, therefore the answer is $\boxed{\frac{1}{4} \text{(C)}}$.

 



\textbf{Solution 4 (It is well known)}

This problem reduces to picking three vectors on the unit circle and asking when the sum of these vectors also lies in the unit circle. But it's well known that the sum of these three vectors is the orthocenter of the triangle whose vertices are the heads of the three vectors. Also, it's well known that the orthocenter of a triangle lies inside its circumcircle iff the triangle is acute. Therefore, this problem is equivalent to asking when three randomly chosen points on a circle form an acute triangle. This is a classic problem, and it's well known that the answer is $\boxed{\frac{1}{4} \text{(C)}}$.

 


	\item (\textit{AMC 12 12B 2010}) A high school basketball game between the Raiders and Wildcats was tied at the end of the first quarter. The number of points scored by the Raiders in each of the four quarters formed an increasing geometric sequence, and the number of points scored by the Wildcats in each of the four quarters formed an increasing arithmetic sequence. At the end of the fourth quarter, the Raiders had won by one point. Neither team scored more than $100$ points. What was the total number of points scored by the two teams in the first half?


 $\textbf{(A)}\ 30 \qquad \textbf{(B)}\ 31 \qquad \textbf{(C)}\ 32 \qquad \textbf{(D)}\ 33 \qquad \textbf{(E)}\ 34$

 \textbf{Answer: }E



\textbf{Solution 1}

Let $a,ar,ar^{2},ar^{3}$ be the quarterly scores for the Raiders. We know $r > 1$ because the sequence is said to be increasing. We also know that each of $a, ar, ar^2, ar^3$ is an integer. We start by showing that \textbf{$r$ must also be an integer.}

 Suppose not, and say $r = m/n$ where $m>n>1$, and $\gcd(m,n)=1$. Then $n, n^2, n^3$ must all divide $a$ so $a=n^3k$ for some integer $k$. Then $S_R = n^3k + n^2mk + nm^2k + m^3k < 100$ and we see that even if $k=1$ and $n=2$, we get $m < 4$, which means that the only option for $r$ is $r=3/2$. A quick check shows that even this doesn't work. Thus $r$ must be an integer.

 Let $a, a+d, a+2d, a+3d$ be the quarterly scores for the Wildcats. Let $S_W = a+(a+d) + (a+2d)+(a+3d) = 4a+6d$. Let $S_R = a+ar+ar^2+ar^3 = a(1+r)(1+r^2)$. Then $S_R<100$ implies that $r<5$, so $r\in \{2, 3, 4\}$. The Raiders win by one point, so
 \[a(1+r)(1+r^2) = 4a+6d+1.\]

            \begin{itemize}
        \item If $r=4$ we get $85a = 4a+6d+1$ which means $3(27a-2d) = 1$, which is not possible with the given conditions.
        \item If $r=3$ we get $40a = 4a+6d+1$ which means $6(6a-d) = 1$, which is also not possible with the given conditions.
        \item If $r=2$ we get $15a = 4a+6d+1$ which means $11a-6d = 1$. Reducing modulo 6 we get $a \equiv 5\pmod{6}$. Since $15a<100$ we get $a<7$. Thus $a=5$. It then follows that $d=9$.
            \end{itemize}
            Then the quarterly scores for the Raiders are $5, 10, 20, 40$, and those for the Wildcats are $5, 14, 23, 32$. Also $S_R = 75 = S_W + 1$. The total number of points scored by the two teams in the first half is $5+10+5+14=\boxed{\textbf{(E)}\ 34}$.

 Note if you don't realize while taking the test that $r$ might not be an integer: since an answer is achieved through casework on the integer value of $r$ and since there is only one right answer, the proof of $r$ being an integer can be skipped on the test (it takes up time).

 



\textbf{Solution 2}

Let $a,ar,ar^{2},ar^{3}$ be the quarterly scores for the Raiders. We know that the Raiders and Wildcats both scored the same number of points in the first quarter so let $a,a+d,a+2d,a+3d$ be the quarterly scores for the Wildcats. The sum of the Raiders scores is $a(1+r+r^{2}+r^{3})$ and the sum of the Wildcats scores is $4a+6d$. Now we can narrow our search for the values of $a,d$, and $r$. Because points are always measured in positive integers, we can conclude that $a$ and $d$ are positive integers. We can also conclude that $r$ is a positive integer by writing down the equation:

 
 \[a(1+r+r^{2}+r^{3})=4a+6d+1\]

 Now we can start trying out some values of $r$. We try $r=2$, which gives

 
 \[15a=4a+6d+1\]

 \[11a=6d+1\]

 We need the smallest multiple of $11$ (to satisfy the <100 condition) that is $\equiv 1 \pmod{6}$. We see that this is $55$, and therefore $a=5$ and $d=9$.

 So the Raiders' first two scores were $5$ and $10$ and the Wildcats' first two scores were $5$ and $14$. 

 
 \[5+10+5+14=34 \longrightarrow \boxed{\textbf{(E)}}\]




\textbf{Solution 3 (Quick Solve?)}

It should become apparent that the geometric ratio is less than or equal to $2$ (try the first quarter $1$, ratio $3$). Obviously, the ratio is more than $1$, so we try the ratio of $2$
$1, 2, 4, 8$ - The Raiders 

 $1, ?, ?, ?$ - The Wildcats

 To find the $?$s, we need the sum of the raiders, which is $15$. $15-1=14$, and $\frac{14}2=7$. So the last term is $7-1=6$. But, now we can't have integrated 2nd and 3rd quarter scores, so we stop.

 $2, 4, 8, 16$ - The Raiders, this also fails (try it out)

 The same goes for $3, 6, 12, 24$, and $4, 8, 16, 32$. With $5, 10, 20, 40$ we get:

 $5, 10, 20, 40$ - The Raiders

 $5, ?, ?, ?$ - The wildcats

 $5+10+20+40=75$, $75-1=74$, $\frac{74}2=37$, $37-5=32$
$5, ?, ?, 32$ - The Wildcats.

 Now, we can finish the arithmetic progression:

 $5, 14, 23, 32$ - The Wildcats

 Thus our answer is $5+10+5+14=34$




\textbf{Solution 4 (solution 2 variant)}

Raiders: $a, am, am^2, am^3$ 

 Wildcats: $a, a + n, a + 2n, a + 3n,$ 

 where $a$ represents both teams' score in the first quarter. Let the Raiders finish with a score of $r$. We know $4a + 6n + 1 = a(1 + m + m^2 + m^3) = r.$ Since $r \le 100$, $1 + m + m^2 + m^3$ can either be $4, 15, 40,$ or $85$. Solve for $a$ in both equations in terms of $r$ and set them equal to each other:

 $\frac{r - (6n + 1)}{4} = \frac{r}{1 + m + m^2 + m^3}.$ We know $1 + m + m^2 + m^3$ can't be $4$ because then $n$ would be a fraction. So let's try $15$:

 $\frac{r - (6n + 1)}{4} = \frac{r}{15} \implies 11r = 90n + 15.$ We know $r$ has to be a multiple of 15 for $a$ to be an integer. After some guess and check with the first few multiples of 15, we find that if $r = 75, n = 9.$ Then $a = 5$. Since this whole case is based on $1 + m + m^2 + m^3 = 15, m = 2$. So our answer is $3a + am + n = \boxed{34}.$

 \textbf{ᴠɪᴅᴇᴏ ѕᴏʟᴜᴛɪᴏɴ ʙʏ ᴘɪ ᴀᴄᴀᴅᴇᴍʏ} 


	\item (\textit{AMC 12 12B 2010}) A geometric sequence $(a_n)$ has $a_1=\sin x$, $a_2=\cos x$, and $a_3= \tan x$ for some real number $x$. For what value of $n$ does $a_n=1+\cos x$?


 

$\textbf{(A)}\ 4 \qquad \textbf{(B)}\ 5 \qquad \textbf{(C)}\ 6 \qquad \textbf{(D)}\ 7 \qquad \textbf{(E)}\ 8$

 \textbf{Answer: }E



\textbf{Solution}

By the defintion of a geometric sequence, we have $\cos^2x=\sin x \tan x$. Since $\tan x=\frac{\sin x}{\cos x}$, we can rewrite this as $\cos^3x=\sin^2x$. 

 The common ratio of the sequence is $\frac{\cos x}{\sin x}$, so we can write

 
 \[a_1= \sin x\]
 \[a_2= \cos x\]
 \[a_3= \frac{\cos^2x}{\sin x}\]
 \[a_4=\frac{\cos^3x}{\sin^2x}=1\]
 \[a_5=\frac{\cos x}{\sin x}\]
 \[a_6=\frac{\cos^2x}{\sin^2x}\]
 \[a_7=\frac{\cos^3x}{\sin^3x}=\frac{1}{\sin x}\]
 \[a_8=\frac{\cos x}{\sin^2 x}=\frac{1}{\cos^2 x}\]

 
Since $\cos^3x=\sin^2x=1-\cos^2x$, we have $\cos^3x+\cos^2x=1 \implies \cos^2x(\cos x+1)=1 \implies \cos x+1=\frac{1}{\cos^2 x}$, which is $a_8$ , making our answer $8 \Rightarrow \boxed{E}$.

 



\textbf{Solution 2}

Notice that the common ratio is $r=\frac{\cos(x)}{\sin(x)}$; multiplying it to $\tan(x)=\frac{\sin(x)}{\cos(x)}$ gives $a_4=1$. Then, working backwards we have $a_3=\frac{1}{r}$, $a_2=\frac{1}{r^2}$ and $a_1=\frac{1}{r^3}$. Now notice that since $a_1=\sin(x)$ and $a_2=\cos(x)$, we need $a_1^2+a_2^2=1$, so $\frac{1}{r^6}+\frac{1}{r^4}=\frac{r^2+1}{r^6}=1\implies r^2+1=r^6$. Dividing both sides by $r^2$ gives $1+\frac{1}{r^2}=r^4$, which the left side is equal to $1+\cos(x)$; we see as well that the right hand side is equal to $a_8$ given $a_4=1$, so the answer is $\boxed{E}$. - mathleticguyyy

 



\textbf{Solution 3}

Note that the common ratio is $\frac{\cos{x}}{\sin{x}} = \cot{x}$. This implies that $\cot{x} \cdot \cos{x} = \tan{x}$. This means $\cos{x} = \frac{\tan{x}}{\cot{x}} = \frac{sin^2{x}}{cos^2{x}} = \frac{1- \cos^2{x}}{\cos^2{x}} = \frac{1}{\cos^2{x}}-1 \implies \cos^3{x} = 1 - \cos^2{x}$. This was a lot of algebra, so take a second to convince yourself that it's all accurate. 

 Next, notice that $a_4 = \tan{x} \cdot \cot{x} = 1$, so $a_{4+i} = \cot^{i}{x}$ by induction. Thus, $a_6 = \cot^2{x} = \frac{\cos^2{x}}{\sin^2{x}} = \frac{\cos^2{x}}{1-\cos^2{x}}$. Substitute our value for $1 - \cos^2{x} = \cos^3{x}$ to get $\cot^2{x} = \sec{x}$. 

 From here, look back at our derivation $\cos^3{x} = 1 - \cos^2{x}$. Rearrange to get $\cos^2{x} \left( \cos{x}+1 \right) = 1 \implies  1+\cos{x} = \sec^2{x} = \cot^4{x}$ by our second derivation. $cot^4{x} = a_8$ which we showed with elementary induction that is left as an exercise for the reader, and our final answer is $\boxed{E}$

 - SumDumNerd

 


	\item (\textit{AMC 12 12B 2010}) Let $a > 0$, and let $P(x)$ be a polynomial with integer coefficients such that


 
$P(1) = P(3) = P(5) = P(7) = a$, and

$P(2) = P(4) = P(6) = P(8) = -a$.


 


 What is the smallest possible value of $a$?


 $\textbf{(A)}\ 105 \qquad \textbf{(B)}\ 315 \qquad \textbf{(C)}\ 945 \qquad \textbf{(D)}\ 7! \qquad \textbf{(E)}\ 8!$

 \textbf{Answer: }B



\textbf{Solution 1}

We observe that because $P(1) = P(3) = P(5) = P(7) = a$, if we define a new polynomial $R(x)$ such that $R(x) = P(x) - a$, $R(x)$ has roots when $P(x) = a$; namely, when $x=1,3,5,7$.

 Thus since $R(x)$ has roots when $x=1,3,5,7$, we can factor the product $(x-1)(x-3)(x-5)(x-7)$ out of $R(x)$ to obtain a new polynomial $Q(x)$ such that $(x-1)(x-3)(x-5)(x-7)(Q(x)) = R(x) = P(x) - a$.

 Then, noting that $P(2)-a=-a-a=-2a$ etc., and plugging in values of $2,4,6,8,$ we get 

 
 \[P(2)-a=(2-1)(2-3)(2-5)(2-7)Q(2) = -15Q(2) = -2a\]
 \[P(4)-a=(4-1)(4-3)(4-5)(4-7)Q(4) = 9Q(4) = -2a\]
 \[P(6)-a=(6-1)(6-3)(6-5)(6-7)Q(6) = -15Q(6) = -2a\]
 \[P(8)-a=(8-1)(8-3)(8-5)(8-7)Q(8) = 105Q(8) = -2a\]
$-2a=-15Q(2)=9Q(4)=-15Q(6)=105Q(8).$Thus, the least value of $a$ must be the $\text{lcm}(15,9,15,105)$.Solving, we receive $315$, so our answer is $\boxed{\textbf{(B)}\ 315}$.

 To complete the solution, we can let $a = 315$, and then try to find $Q(x)$. We know from the above calculation that $Q(2)=42, Q(4)=-70, Q(6)=42$, and $Q(8)=-6$. Then we can let $Q(x) = T(x)(x-2)(x-6)+42$, getting $T(4)=28, T(8)=-4$. Let $T(x)=L(x)(x-8)-4$, then $L(4)=-8$. Therefore, it is possible to choose $T(x) = -8(x-8)-4 = -8x + 60$, so the goal is accomplished. As a reference, the polynomial we get is

 
 \[P(x) = (x-1)(x-3)(x-5)(x-7)((-8x + 60)(x-2)(x-6)+42) + 315\]
 \[= -8 x^7+252 x^6-3248 x^5+22050 x^4-84392 x^3+179928 x^2-194592 x+80325\]




\textbf{Solution 2}

The evenly-spaced data suggests using \href{https://artofproblemsolving.com/wiki/index.php?title=Discrete\_derivative}{discrete derivatives} to tackle this problem.  First, note that any polynomial of degree $n$
$P(x) = a_0 + a_1 x + a_2 x^2 + \ldots + a_n x^n$

 can also be written as

 $P(x) = b_0 + b_1 (x-1) + b_2 (x-1)(x-2) + \ldots + b_n (x-1)(x-1) \cdots (x-n)$.

 Moreover, the coefficients $a_i$ are integers for $i=1, 2, \ldots n$ iff the coefficients $b_i$ are integers for $i=1, 2, \ldots n$.  This latter form is convenient for calculating discrete derivatives of $P(x)$.

 The discrete derivative of a function $f(x)$ is the related function $\Delta f(x)$ defined as

 $\delta f(x) = f(x+1) - f(x)$.

 With this definition, it's easy to see that for any positive integer $k$ we have

 $\Delta [(x-1)(x-2)\cdots(x-k)] = k(x-1)(x-2)\cdots(x-[k-1])$.

 This in turn allows us to use successive discrete derivatives evaluated at $x=1$ to calculate all of the coefficients $b_i$ using

 $P(1) = b_0$, $\Delta P(1) = b_1$, $\Delta^2 P(1) = 2 b_2$, $\ldots$, $\Delta^7 P(1) = 7! b_7$.

 We can also calculate the following table of discrete derivatives based on the data points given in the problem statement:

 <table cellpadding="3" frame="box" rules="all"><tbody><tr><th></th><th colspan="8">$x$</th></tr><tr><td align="right"></td><td align="center">$1$</td><td align="center">$2$</td><td align="center">$3$</td><td align="center">$4$</td><td align="center">$5$</td><td align="center">$6$</td><td align="center">$7$</td><td align="center">$8$</td></tr><tr><td align="right">$P(x)$</td><td align="right">$a$</td><td align="right">$-a$</td><td align="right">$a$</td><td align="right">$-a$</td><td align="right">$a$</td><td align="right">$-a$</td><td align="right">$a$</td><td align="right">$-a$</td></tr><tr><td align="right">$\Delta P(x)$</td><td align="right">$-2a$</td><td align="right">$2a$</td><td align="right">$-2a$</td><td align="right">$2a$</td><td align="right">$-2a$</td><td align="right">$2a$</td><td align="right">$-2a$</td><td></td></tr><tr><td align="right">$\Delta^2 P(x)$</td><td align="right">$4a$</td><td align="right">$-4a$</td><td align="right">$4a$</td><td align="right">$-4a$</td><td align="right">$4a$</td><td align="right">$-4a$</td><td></td><td></td></tr><tr><td align="center" colspan="9">$\vdots$</td></tr><tr><td align="right">$\Delta^7 P(x)$</td><td align="right">$-2^7 a$</td><td></td><td></td><td></td><td></td><td></td><td></td><td></td></tr></tbody></table>

 Thus we can read down the column for $x=1$ to find that $k! b_k = (-2)^k a$ for $k = 0, 1, \ldots, 7$.  Interestingly, even if we choose $P(x)$ to have degree greater than $7$, the $8$ coefficients of lowest order always satisfy these conditions.  Moreover, it's straightforward to show that the $P(x)$ of degree $7$ with $b_k$ satisfying these conditions will fit the data given in the problem statement.  Thus, to find minimal necessary and sufficient conditions on the value of $a$, we need only consider these $8$ equations.  As a result, $P(x)$ with integer coefficients fitting the given data exists iff $k!$ divides $2^k a$ for $k = 0, 1, \ldots, 7$.  In other words, it's necessary and sufficient that

 
 \[0! | a\],

 
 \[1! | 2a\],

 
 \[2! | 2^2 a\],

 
 \[3! | 2^3 a\],

 
 \[4! | 2^4 a\],

 
 \[5! | 2^5 a\],

 
 \[6! | 2^6 a, \text{and}\]

 \[7! | 2^7 a.\]

 The last condition holds if $7 \cdot 3 \cdot 5 \cdot 3 = 315$ divides evenly into $a$.  Since such $a$ will also satisfy the first $7$ conditions, our answer is $\boxed{\textbf{(B)}\ 315}$.

 



\textbf{Solution 3 (2 liner)}

We have $P(x)-a=(x-1)(x-3)(x-5)(x-7)Q(x).$ Plugging in $x=2,4,6,$ and $8$, we get
 \[Q(2)=\dfrac{2}{15}a\]
 \[Q(4)=-\dfrac{2}{9}a\]
 \[Q(6)=-\dfrac{2}{15}a\]
 \[Q(8)=\dfrac{2}{105}a.\]Thus, $a=\text{lcm}(9,15,105)=\boxed{\textbf{(B)}\ 315}.$

 -jb2015007

 


	\item (\textit{AMC 12 12B 2010}) Let $ABCD$ be a cyclic quadrilateral. The side lengths of $ABCD$ are distinct integers less than $15$ such that $BC\cdot CD=AB\cdot DA$. What is the largest possible value of $BD$?


 $\textbf{(A)}\ \sqrt{\dfrac{325}{2}} \qquad \textbf{(B)}\ \sqrt{185} \qquad \textbf{(C)}\ \sqrt{\dfrac{389}{2}} \qquad \textbf{(D)}\ \sqrt{\dfrac{425}{2}} \qquad \textbf{(E)}\ \sqrt{\dfrac{533}{2}}$

 \textbf{Answer: }D



\textbf{Solution 1}

Let $AB = a$, $BC = b$, $CD = c$, and $AD = d$.  We see that by the Law of Cosines on $\triangle ABD$ and $\triangle CBD$, we have:

 $BD^2 = a^2 + d^2 - 2ad\cos{\angle BAD}$.  

 $BD^2 = b^2 + c^2 - 2bc\cos{\angle BCD}$.

 We are given that $ad = bc$ and $ABCD$ is a cyclic quadrilateral. As a property of cyclic quadrilaterals, opposite angles are supplementary so $\angle BAD = 180 - \angle BCD$, therefore $\cos{\angle BAD} = -\cos{\angle BCD}$. So, $2ad\cos{\angle BAD} = -2bc\cos{\angle BCD}$. 

 Adding, we get $2BD^2 = a^2 + b^2 + c^2 + d^2$.

 We now look at the equation $ad = bc$.  Suppose that $a = 14$.  Then, we must have either $b$ or $c$ equal $7$.  Suppose that $b = 7$.  We let $d = 6$ and $c = 12$.  

 $2BD^2 = 196 + 49 + 36 + 144 = 425$, so our answer is $\boxed{\textbf{(D)} \sqrt{\frac{425}{2}}}$.

 



\textbf{Solution 2}

As previous, let $AB = a$, $BC = b$, $CD = c$, and $AD = d$ and proceed as above until the addition step. From here, note

 $a^2+d^2-b^2-c^2=4ad\cos{\theta}$
$2ad\cos{\theta} = \frac{1}{2}(a^2+d^2-b^2-c^2)$

 Therefore, $a^2+d^2-2ad\cos{\theta} = \frac{1}{2}(a^2+b^2+c^2+d^2)$. Note that because of the bounds of $\cos{\theta}$, this is in the interval $[-4ad, 4ad]$. We want to maximize this, so we pick the largest possible side length, $a=14$. Then, we must have $b=7$ to satisfy $ad=bc$. We cannot pick $13$ because it is prime and the integers are distinct, so we choose $c=12$ and $d=6$.This is the optimal set of side lengths that maximizes our expression, which comes out to $\frac{1}{2}(196+49+144+36) = \frac{425}{2}$. This is $BD^2$, so the answer is the square root of it, which is $\boxed{\textbf{(D)} \sqrt{\frac{425}{2}}}$

 - SumDumNerd

 


	\item (\textit{AMC 12 12B 2010}) Monic quadratic polynomials $P(x)$ and $Q(x)$ have the property that $P(Q(x))$ has zeros at $x=-23, -21, -17,$ and $-15$, and $Q(P(x))$ has zeros at $x=-59,-57,-51$ and $-49$. What is the sum of the minimum values of $P(x)$ and $Q(x)$? 


 $\textbf{(A)}\ -100 \qquad \textbf{(B)}\ -82 \qquad \textbf{(C)}\ -73 \qquad \textbf{(D)}\ -64 \qquad \textbf{(E)}\ 0$

 \textbf{Answer: }A



\textbf{Solution 1}

$P(x) = (x - a)^2 - b, Q(x) = (x - c)^2 - d$. Notice that $P(x)$ has roots $a\pm \sqrt {b}$, so that the roots of $P(Q(x))$ are the roots of $Q(x) = a + \sqrt {b}, a - \sqrt {b}$. For each individual equation, the sum of the roots will be $2c$ (symmetry or Vieta's). Thus, we have $4c = - 23 - 21 - 17 - 15$, or $c = - 19$. Doing something similar for $Q(P(x))$ gives us $a = - 54$.We now have $P(x) = (x + 54)^2 - b, Q(x) = (x + 19)^2 - d$. Since $Q$ is monic, the roots of $Q(x) = a + \sqrt {b}$ are "farther" from the axis of symmetry than the roots of $Q(x) = a - \sqrt {b}$. Thus, we have $Q( - 23) = - 54 + \sqrt {b}, Q( -21) =- 54 - \sqrt {b}$, or $16 - d = - 54 + \sqrt {b}, 4 - d = - 54 - \sqrt {b}$. Adding these gives us $20 - 2d = - 108$, or $d = 64$. Plugging this into $16 - d = - 54 + \sqrt {b}$, we get $b = 36$.The minimum value of $P(x)$ is $- b$, and the minimum value of $Q(x)$ is $- d$. Thus, our answer is $- (b + d) = - 100$, or answer $\boxed{\textbf{(A)}}$.

 


 



\textbf{Solution 2 (Bash)}

Let $P(x) = x^2 + Bx + C$ and $Q(x) = x^2 + Ex + F$.

 Then $P(Q(x))$ is $(x^2 + Ex + F)^2 + B(x^2 + Ex + F) + C$, which simplifies to:

 $P(Q(x)) = x^4 + 2Ex^3 + (E^2 + 2F + B)x^2 + (2EF + BE)x + (F^2 + BF + C)$

 We can find $Q(P(x))$ by simply doing $B\Leftrightarrow E$ and $C \Leftrightarrow F$ to get:

 $Q(P(x)) = x^4 + 2Bx^3 + (B^2 + 2C + E)x^2 + (2BC + BE)x + (C^2 + EC + F)$

 The sum of the zeros of $P(Q(x))$ is $-76$. From Vieta, the sum is $-2E$. Therefore, $E = 38$.

 The sum of the zeros of $Q(P(x))$ is $-216$. From Vieta, the sum is $-2B$. Therefore, $B = 108$.

 Plugging in, we get:

 $P(Q(x)) = x^4 + 76x^3 + (1552 + 2F)x^2 + (76F + 4104)x + (F^2 + 108F + C)$$Q(P(x)) = x^4 + 216x^3 + (11702 + 2C)x^2 + (216C + 4104)x + (C^2 + 38C + F)$

 Let's tackle the $x^2$ coefficients, which is the sum of the six double-products possible. Since $23 \cdot (21 + 17 + 15) + 21 \cdot (17 + 15) + 17 \cdot 15$ gives the sum of these six double products of the roots of $P(Q(x))$, we have:

 $1552 + 2F = 23 \cdot (21 + 17 + 15) + 21 \cdot (17 + 15) + 17 \cdot 15$
$1552 + 2F = 2146$
$F = 297$

 Similarly with $Q(P(x))$, we get:

 $11702 + 2C = 59(57 + 51 + 49) + 57(51 + 49) + 51(49)$
$11702 + 2C = 17462$
$C = 2880$

 Thus, our polynomials are $P(x) = x^2 + 108x + 2880$ and $Q(x) = x^2 + 38x + 297$.

 The minimum value of $P(x)$ happens at $x = -\frac{108}{2} = -54$, and is $54^2 - 108 \cdot 54 + 2880 = 2880 - 54^2$.

 The minimum value of $Q(x)$ happens at $x = -\frac{38}{2} = -19$, and is $19^2 - 38 \cdot 19 + 297 = 297 - 19^2$.

 The sum of these minimums is $2880 +297 - 54^2 - 19^2 = \boxed{-100}$. -srisainandan6

 



\textbf{Solution 3 (Mild Bash)}

Let $P(x) = x^2 - (a+b)x + ab$ and $Q(x) = x^2 - (c+d)x + cd$. Notice that the roots of $P(x)$ are $a,b$ and the roots of $Q(x)$ are $c,d.$ Then we get:

 
 \begin{align*} P(Q(x)) &= a, b \\ x^2 - (c+d)x + cd &= a, b \end{align*}
The two possible equations are then $x^2 - (c+d)x + cd-a=0$ and $x^2 - (c+d)x + cd-b=0$. The solutions are $-23, -21, -17, -15$. From Vieta's we know that the total sum $2(c+d) = -76 \implies c+d = -38$ so the roots are paired $-23, -15$ and $-21, -17$. Then, $cd - a = 23*15$ and $cd - b = 21*17$.

 We can similarly get that $ab - c = 59*49$ and $ab - d = 57*51$, and $a+b = -108$. Add the first two equations to get 
 \[2cd - (a+b) = 23*15 + 21*17 \implies cd = \frac{23*15+21*17 - 108}{2} = 297.\] This means $Q(x) = x^2 + 38x + 297$.

 Once more, we can similarly obtain 
 \[ab = \frac{59*49 + 57*51 - 38}{2} = 2880.\] Therefore $P(x) = x^2 + 108x + 2880$.

 Now we can find the minimums to be 
 \[19^2 - 19*38 + 297 = -64\] and 
 \[54^2 - 54*108 + 2880 = -36.\] Summing, the answer is $\boxed{\textbf{(A)} -100}.$




\textbf{Solution 4}

Let $P(x) = (x+a)(x+b)$, $Q(x) = (x+c)(x+d)$. 

 $P(Q(x)) = (x^2 + cx + dx + cd + a)(x^2 + cx + dx + cd + b)$
$Q(P(x)) = (x^2 + ax + bx + ab + c)(x^2 + ax + bx + ab + d)$

 Notice how the coefficient for $x$ has to be the same for the two quadratics that are multiplied to create $P(Q(x))$, and $Q(P(x))$.

 $P(Q(x)) = (x+ 23)(x+ 21)(x+ 17)(x+ 15) = (x^2 + 38x + 345)(x^2 + 38x + 357)$
$Q(P(x)) = (x+ 59)(x+ 57)(x+ 51)(x+ 49) = (x^2 + 108x + 2891)(x^2 + 108x + 2907)$
$c + d = 38$, $cd + a = 345$, $cd + b = 357$, $a + b = 108$, $ab + c = 2891$, $ab + d = 2907$
$cd = \frac{345 + 357 - 108}{2} = 297$, $297 = 3^3 \cdot 11$, $c = 27$, $d = 11$
$a = 345 - 27*11 = 48$, $b = 357 - 27*11 = 60$
$P(x) = (x+48)(x+60) = x^2 + 108x + 2880 = (x+54)^2 - 36$
$Q(x) = (x+27)(x+11) = x^2 + 38x + 297 = (x+19)^2 -64$
$-36 -64 = \boxed{\textbf{(A) } -100}$.

 



\textbf{Solution 5}

Because $P$ is a quadratic, then $P$ only has two roots; however, $P(Q(x))$ has four roots, meaning that two of the values of $Q(x)$ must repeat themselves. 

 
Notice that the four roots of $P(Q(x))$ are symmetric around $x=-19$. This means that $Q(x)$'s graph has vertex $-19$; therefore, $Q(-17)=Q(-21$) and $Q(-15)=Q(-23)$. Similarly, we can deduce that $P(x)$ has vertex $-54$ such that $P(-51)=P(-57)$ and $P(-49)=P(-59)$. 

 
If we have $P(x)=x^2+bx+c$, then since the vertex is $-54,\text{ }b=108$; similarly, we get that $Q(x)=x^2+nx+m=x^2+38x+m$. Since the roots of $P(x)$ are $Q(-17)$ and $Q(-15)$ (or $Q(-23)$ and $Q(-21)$, they are all equal), by simplifying we get that the roots of $P(x)$ are $m-357$ and $m-345$. 

 
By Vieta's Formulas, $(m-357)+(m-345)=-108$, therefore $m=297$. Similarly, we can get $c=2880$. By applying the formula for minima, we get that the minimum for $Q(x)$ is $-64$ and the minimum for $P(x)$ is $-36$, and therefore the answer is $\boxed{-100}$. $\square$



	\item (\textit{AMC 12 12B 2010}) The set of real numbers $x$ for which 


 
 \[\dfrac{1}{x-2009}+\dfrac{1}{x-2010}+\dfrac{1}{x-2011}\ge1\]


 is the union of intervals of the form $a<x\le b$. What is the sum of the lengths of these intervals?


 $\textbf{(A)}\ \dfrac{1003}{335} \qquad \textbf{(B)}\ \dfrac{1004}{335} \qquad \textbf{(C)}\ 3 \qquad \textbf{(D)}\ \dfrac{403}{134} \qquad \textbf{(E)}\ \dfrac{202}{67}$

 \textbf{Answer: }C



\textbf{Solution 1}

Because the right side of the inequality is a horizontal line, the left side can be translated horizontally by any value and the intervals will remain the same. For simplicity of calculation, we will find the intervals where 
 \[\frac{1}{x+1}+\frac{1}{x}+\frac{1}{x-1}\ge1\]We shall say that $f(x)=\frac{1}{x+1}+\frac{1}{x}+\frac{1}{x-1}$. $f(x)$ has three vertical asymptotes at $x=\{-1,0,1\}$. As the sum of decreasing hyperbolas, the function is decreasing at all intervals. Values immediately to the left of each asymptote approach negative infinity, and values immediately to the right of each asymptote approach positive infinity. In addition, the function has a horizontal asymptote at $y=0$. The function intersects $1$ at some point from $x=-1$ to $x=0$, and at some point from $x=0$ to $x=1$, and at some point to the right of $x=1$. The intervals where the function is greater than $1$ are between the points where the function equals $1$ and the vertical asymptotes.

 If $p$, $q$, and $r$ are values of x where $f(x)=1$, then the sum of the lengths of the intervals is $(p-(-1))+(q-0)+(r-1)=p+q+r$.

 
 \[\frac{1}{x+1}+\frac{1}{x}+\frac{1}{x-1}=1\]
 \[\implies x(x-1)+(x-1)(x+1)+x(x+1)=x(x-1)(x+1)\]
 \[\implies x^3-3x^2-x+1=0\]

 And now our job is simply to find the sum of the roots of $x^3-3x^2-x+1$.Using \href{https://artofproblemsolving.com/wiki/index.php?title=Vieta\%27s\_formulas}{Vieta's formulas}, we find this to be $3$ $\Rightarrow\boxed{C}$.

 \textit{NOTE'}: For the AMC, one may note that the transformed inequality should not yield solutions that involve big numbers like 67 or 134, and immediately choose $C$.

 



\textbf{Solution 2}

As in the first solution, note that the expression can be translated into 
 \[\frac{1}{x+1}+\frac{1}{x}+\frac{1}{x-1}\ge1\] without affecting the interval lengths.

 This simplifies into 
 \[\frac{-x^3+3x^2+x-1}{(x)(x+1)(x-1)}\ge0,\] so 
 \[-x^3+3x^2+x-1\ge0.\]Each interval is $(-1, a), (0, b), (1, c)$, where $a$, $b$, and $c$ are the roots of $-x^3+3x^2+x-1=0$so the total length is $a+b+c$, which is the sum of the roots, or $\boxed3$.

 



\textbf{Solution 3}

Note that all of the answers but $C$ have weird factors of 2010, but 2010 is a random number (set $x'=x-2010$). So therefore the answer is $\fbox{C(3)}$



	\item (\textit{AMC 12 12B 2010}) For every integer $n\ge2$, let $\text{pow}(n)$ be the largest power of the largest prime that divides $n$. For example $\text{pow}(144)=\text{pow}(2^4\cdot3^2)=3^2$. What is the largest integer $m$ such that $2010^m$ divides


 
$\prod_{n=2}^{5300}\text{pow}(n)$?


 


 

$\textbf{(A)}\ 74 \qquad \textbf{(B)}\ 75 \qquad \textbf{(C)}\ 76 \qquad \textbf{(D)}\ 77 \qquad \textbf{(E)}\ 78$

 \textbf{Answer: }D



\textbf{Solution}

Because 67 is the largest prime factor of 2010, it means that in the prime factorization of $\prod_{n=2}^{5300}\text{pow}(n)$, there'll be $p_1 ^{e_1} \cdot p_2 ^{e_2} \cdot .... 67^x ...$ where $x$ is the desired value we are looking for. Thus, to find this answer, we need to look for the number of times $67$ is incorporated into the giant product.

 All numbers $n=67 \cdot x$, given $x = p_1 ^ {e_1} \cdot p_2 ^{e_2} \cdot ... \cdot p_m ^ {e_m}$ such that for any integer $x$ between $1$ and $m$, prime $p_x$ must be less than $67$, contributes a 67 to the product. Considering $67 \cdot 79 < 5300 < 67 \cdot 80$, the possible values of x are $1,2,...,70,72,74,...78$, since $x=71,73,79$ are primes that are greater than 67. However, $\text{pow}\left(67^2\right)$ contributes two $67$s to the product, so we must count it twice. Therefore, the answer is $70 + 1 + 6 = \boxed{77} \Rightarrow \boxed{D}$.

 \textbf{Similar Solution} After finding the prime factorization of $2010=2\cdot3\cdot5\cdot67$, divide $5300$ by $67$ and add $5300$ divided by $67^2$ in order to find the total number of multiples of $67$ between $2$ and $5300$. $\lfloor\frac{5300}{67}\rfloor+\lfloor\frac{5300}{67^2}\rfloor=80$ Since $71$,$73$, and $79$ are prime numbers greater than $67$ and less than or equal to $80$, subtract $3$ from $80$ to get the answer $80-3=\boxed{77}\Rightarrow\boxed{D}$.

 



\textbf{Solution 3 (quick, observations, similar to 2)}

$2010 = 2\cdot3\cdot5\cdot67$.

 There will probably be enough $2$'s, $3$'s, and $5$'s to not count them. Therefore we must count the number of $67$'s in $\prod_{n=2}^{5300}\text{pow}(n)$.

 Notice $\text{pow}\left(n\right)=67$ if and only if $n={67}\cdot{b}$, where $\text{pow}\left(b\right)$ is less than OR EQUAL TO $67$. Therefore, we have $b\le\lfloor{\frac{5300}{67}}\rfloor$, or $b\le{79}$. Therefore we have $79$ options plus one ($67^2$ will add 2 $67$'s) minus $3$ (for $71$, $73$, and $79$) to get $\boxed{77}\Rightarrow\boxed{D}$.

 \textbf{Need Discussion + Clarification?} Concern: How do we know that we only have to check 67? There is no solid relationship between 67 being the largest prime factor in 2010 and 67 giving the smallest result of 77.

 
Response to concern: Ok, 67 is waaaay larger than any of the other prime factors of 2010. 2001's prime factors are close so there might be some ambiguity but if you can't see why given intervals of 67, 2s 3s and 5s don't need to be accounted for, you probably can't solve this problem .-.. Details in the Discussions page of this Article: \href{https://artofproblemsolving.comhttp://artofproblemsolving.com/wiki/index.php?title=Talk:2010\_AMC\_12B\_Problems/Problem\_25}{http://artofproblemsolving.com/wiki/index.php?title=Talk:2010\_AMC\_12B\_Problems/Problem\_25}

 
Response to concern:When I solved it, I actually ended up trying 2, 3, and 5 just to be safe. It actually ends up not being that bad (since for each one, we only need to consider numbers that only have that prime and lower primes in its prime factorization, so you can just divide 5300 by your prime, find how many numbers less than that new number are only divisible by lower primes, divide by your prime again, and repeat). All said and done, it only ended up taking about 7-8 minutes of mental math (ran out of paper oops), and I think that reasonably speaking, many who are good enough to have a decent chance of solving this 25 can check 2, 3, and 5 on paper in 5 minutes.

 


\end{enumerate}